\documentclass{report}
\usepackage{amsfonts, amsmath}
\newcommand\R{{\mathbb R}}
\newcommand\Q{{\mathbb Q}}
\newcommand\N{{\mathbb N}}
\pagestyle{empty}
\begin{document}
\large
\centerline{\Huge \bf Analisi Matematica I}
\vskip.2cm
\centerline{\Huge Primo Esonero}
\vskip.1cm
\centerline{\huge \tt A} 
\renewcommand{\thefootnote}{ } 
\footnote{\sl Avete 3 ore di tempo. Ogni esercizio
vale sei punti. La sufficienza si ottiene con un punteggio $\geq$
18. {\bf Solo le risposte chiaramente giustificate saranno prese in
considerazione.} {\scshape Elaborati illeggibili o confusi verranno ignorati.}
{\sffamily La
sintesi sar\`a particolarmente apprezzata.}}
\renewcommand{\thefootnote}{\arabic{footnote}} 

\vskip1cm
\begin{enumerate}
\item Si verifichi che la funzione $f(x)=e^x+1$ \`e monotona crescente
(cio\`e se $x\geq y$ allora $f(x)\geq f(y)$).

\item Si calcoli il limite
\[
\lim_{n\to\infty} n\left\{\left(1+\frac 1n\right)^{\frac
14}-1\right\}.
\]

\item Si studi la convergenza della serie
\[
\sum_{n=1}^\infty n! 2^{-n^2}\, .
\]
\item Si consideri la funzione 
\[
f(x)=\frac 12-e^{-x^2}.
\]
Si dimostri che l'equazione $f(x)=0$ ha almeno due soluzioni in $\R$.

\item Si consideri la funzione
\[
f(x)= e^{\frac 12 \sqrt{x}+3}\quad\text{per }x\in\R^+.
\]
Si determini $f^{-1}$.
\item Sia $a=0.\overline{123}$ dove la barra sopra $123$ significa che
le cifre sono ripetute indefinitivamente (cio\`e
$a=0.123123123\dots$). Si dimostri che $a\in\Q$.

\end{enumerate}
\end{document}
